\section{The complete positive-phase affine table}\label{app:affine}

The table uses the rectangle convention \eqref{eq:rectangle}.
Each row asserts an affine identity $g^\ell(x,y)=(x,y)$ on
its entire residue cell; its last column lists every point in
that cell having proper-divisor period. A dash means there is
no such point. The verification procedure in
Proposition~\ref{prop:core} requires only the six values in
\eqref{eq:sigma}, affine substitution through at most $20$ steps,
and two-variable linear equations. Thus the table is a finite
proof table included in the article, not a reference to an
unprovided degree census.

{\small
\begin{longtable}{ccccc}
\caption{All $36$ positive-phase core cells and their central exceptions.}
\label{tab:affine}\\
\toprule
$(e,f)$ & $\ell$ & $(L_x,U_x)$ & $(L_y,U_y)$ &
Exceptional point: period\\
\midrule
\endfirsthead
\multicolumn{5}{c}{Table \thetable\ continued}\\
\toprule
$(e,f)$ & $\ell$ & $(L_x,U_x)$ & $(L_y,U_y)$ &
Exceptional point: period\\
\midrule
\endhead
\midrule
\multicolumn{5}{r}{Continued on next page}\\
\endfoot
\bottomrule
\endlastfoot
$(0,0)$ & $4$ & $(0,0)$ & $(0,0)$ & $(0,0):1$\\
$(0,1)$ & $12$ & $(1,-1)$ & $(2,0)$ & $(0,1):6$\\
$(0,2)$ & $20$ & $(2,-2)$ & $(4,0)$ & $(0,2):5$\\
$(0,3)$ & $4$ & $(0,0)$ & $(0,0)$ & ---\\
$(0,4)$ & $20$ & $(2,-2)$ & $(0,-4)$ & $(0,-2):5$\\
$(0,5)$ & $12$ & $(1,-1)$ & $(0,-2)$ & $(0,-1):6$\\
\midrule
$(1,0)$ & $12$ & $(2,0)$ & $(1,-1)$ & $(1,0):6$\\
$(1,1)$ & $12$ & $(2,0)$ & $(2,0)$ & $(1,1):6$\\
$(1,2)$ & $20$ & $(3,-1)$ & $(4,0)$ & $(1,2):5$\\
$(1,3)$ & $20$ & $(0,-4)$ & $(2,-2)$ & ---\\
$(1,4)$ & $20$ & $(0,-4)$ & $(3,-1)$ & ---\\
$(1,5)$ & $20$ & $(3,-1)$ & $(1,-3)$ & $(1,-1):5$\\
\midrule
$(2,0)$ & $20$ & $(4,0)$ & $(2,-2)$ & $(2,0):5$\\
$(2,1)$ & $20$ & $(4,0)$ & $(3,-1)$ & $(2,1):5$\\
$(2,2)$ & $20$ & $(1,-3)$ & $(1,-3)$ & ---\\
$(2,3)$ & $12$ & $(0,-2)$ & $(1,-1)$ & ---\\
$(2,4)$ & $12$ & $(0,-2)$ & $(2,0)$ & ---\\
$(2,5)$ & $20$ & $(1,-3)$ & $(4,0)$ & ---\\
\midrule
$(3,0)$ & $4$ & $(0,0)$ & $(0,0)$ & ---\\
$(3,1)$ & $20$ & $(2,-2)$ & $(0,-4)$ & ---\\
$(3,2)$ & $12$ & $(1,-1)$ & $(0,-2)$ & ---\\
$(3,3)$ & $4$ & $(0,0)$ & $(0,0)$ & ---\\
$(3,4)$ & $12$ & $(1,-1)$ & $(2,0)$ & ---\\
$(3,5)$ & $20$ & $(2,-2)$ & $(4,0)$ & ---\\
\midrule
$(4,0)$ & $20$ & $(0,-4)$ & $(2,-2)$ & $(-2,0):5$\\
$(4,1)$ & $20$ & $(3,-1)$ & $(0,-4)$ & ---\\
$(4,2)$ & $12$ & $(2,0)$ & $(0,-2)$ & ---\\
$(4,3)$ & $12$ & $(2,0)$ & $(1,-1)$ & ---\\
$(4,4)$ & $20$ & $(3,-1)$ & $(3,-1)$ & ---\\
$(4,5)$ & $20$ & $(0,-4)$ & $(1,-3)$ & $(-2,-1):5$\\
\midrule
$(5,0)$ & $12$ & $(0,-2)$ & $(1,-1)$ & $(-1,0):6$\\
$(5,1)$ & $20$ & $(1,-3)$ & $(3,-1)$ & $(-1,1):5$\\
$(5,2)$ & $20$ & $(4,0)$ & $(1,-3)$ & ---\\
$(5,3)$ & $20$ & $(4,0)$ & $(2,-2)$ & ---\\
$(5,4)$ & $20$ & $(1,-3)$ & $(0,-4)$ & $(-1,-2):5$\\
$(5,5)$ & $12$ & $(0,-2)$ & $(0,-2)$ & $(-1,-1):6$\\
\end{longtable}
}

